The section now starts from the abstract definition of Haar measure and obtains hyperbolic area as the measure induced on $G/K$. Its companion partially formalizes the abstract results: \texttt{haar\_unique\_scale} proves uniqueness for two left Haar measures, and \texttt{haar\_inverse} proves that inversion takes a nonzero regular left-invariant measure to a nonzero regular right-invariant measure. The right-Haar uniqueness statement follows in the prose by inversion but does not yet have a separate exported proof here. Steven's correspondence review remains pending.

The shared module \texttt{TraceFormulas.MeasureTheory.HaarBasics} uses Mathlib's genuine Borel measures and topological groups. The companion starts from regularity, left invariance and nonzeroness; it derives the library's positivity-on-open-sets condition rather than assuming a different meaning of Haar measure. The unique scalar has type \texttt{NNReal}, is strictly positive, and therefore is a finite positive real number. The formal uniqueness proof uses Mathlib's general invariant-measure theorem and normalization on a compact set with nonempty interior; the appendix instead gives a compact-support approximation and Fubini argument. Appendix A.2's comparison records these details and the regularity convention.

The former real-line lemma and its unchanged mathematical proofs now belong to \hyperref[measure.translation-invariant-line]{Appendix A.1}, with a companion at that location. They no longer stand in for this section's general Haar statements.

The Iwasawa product measure, its left invariance, the quotient pushforward to $\mathbb H$, and the lift of an arbitrary invariant Radon measure remain formalization targets. The prose now deduces invariant-area uniqueness from Haar uniqueness; it no longer uses the horizontal-slice and logarithmic-coordinate proof. The checked abstract uniqueness theorem alone does not discharge the construction and measure-transport obligations.

The diagonal-factor step in the \hyperref[modular.unimodular]{unimodularity proof} now has three companion declarations. The declaration \texttt{quarter\_turn\_conjugation} proves $k(\pi/2)a(y)k(\pi/2)^{-1}=a(y)^{-1}$ for the existing determinant-one matrices, with $y>0$. The declaration \texttt{positive\_reciprocal\_fixed\_point} proves that a positive real number equal to its inverse is $1$, by deriving its square is $1$ and ruling out $-1$. Finally, \texttt{positive\_character\_on\_dilation} proves that any homomorphism from $G$ to the nonzero real numbers, with an explicit positivity hypothesis, takes value $1$ on every $a(y)$. The target is Mathlib's group of real units, so these are genuine nonzero real values. The shared module \texttt{TraceFormulas.Geometry.IwasawaCharacters} proves the conjugation calculation and the more general fact about an element conjugate to its inverse. This matches the edited algebraic argument. Applying it to the actual modular function, proving that function's homomorphism and continuity, completing the $N$ and $K$ factors and the passage to $\overline G$, and proving the central pushforward remain open in Lean.

The \hyperref[modular.quotient-measure]{geometric quotient definition} uses the same measure on the canonically identified orbit spaces $\Gamma\backslash\mathbb H$ and $\overline\Gamma\backslash\mathbb H$. The earlier checked action-kernel calculation establishes that $-I$ acts trivially, and Mathlib's \texttt{UpperHalfPlane.SL\_neg\_smul} identifies the actions of $g$ and $-g$. These facts support the orbit description; they do not yet construct the effective quotient or its measure. The proof of independence of representatives uses the disjoint translates of $\mathcal F^\circ$ indexed by $\overline\Gamma$. Indexing those translates by all of $\Gamma$ would repeat each region twice.

Other targets are genuine geometric quotient measures and unfolding, the finite-index identity, and the fiber factor $m$ for the different quotient $\Lambda\backslash G$. The theorem now defines $m=2$ if $-I\in\Lambda$ and $m=1$ otherwise; this is the same quantity $|\Lambda\cap\{I,-I\}|$ used before. It counts the elements of each fiber of the projection $\Lambda\to\overline\Lambda$. The projection is an isomorphism when $m=1$ and is two-to-one when $m=2$. This subgroup projection and the resulting covolume formulas remain formalization targets. Compact fibers carry probability measures and discrete subgroups carry counting measure. The normalization is $\nu(Eh)=\Delta_L(h)\nu(E)$, so right translation in an integral contributes the reciprocal. No constant defined to equal the intended area or covolume is a proof of its calculation.

The expanded proof gives an explicit roadmap for the three remaining volume formulas: assemble a measurable representative region from coset translates and sum their areas; construct matrix-group representatives with $0\leq\theta<2\pi/m$ and integrate $d\theta/(2\pi)$; then construct effective-group representatives with $0\leq\theta<\pi$ and integrate $d\theta/\pi$. A formal proof must justify these representative sets and their null exceptions for the actual quotient measures before evaluating the integrals. The existing companion does not yet prove these steps.
